\documentclass[10pt,a4paper]{amsart}
\usepackage[margin=19mm]{geometry}
\usepackage{amsmath,amssymb,mathtools}
\usepackage[scr=rsfs,cal=euler]{mathalfa}
\usepackage{xfrac}
\usepackage[unicode,hidelinks,bookmarks=false]{hyperref}
\newcommand{\ie}{i.e.\ }
\newcommand{\cf}{cf.\ }
\newcommand{\resp}{respectively\ }
\newcommand{\Subsection}{\subsection}
\providecommand{\nobreakdash}{\nobreak\hbox{-}}
\setlength{\emergencystretch}{2em}
\multlinegap0pt
\usepackage{mathtools}
\usepackage{amssymb}
\usepackage[shortlabels]{enumitem}
\usepackage{xcolor}
\newtheorem{theorem}{Theorem}
\newcommand{\C}{\mathcal C}
\newcommand{\DFrm}{\mathbf{DFrm}}
\newcommand{\Null}{\mathcal N}
\newcommand{\SDFrm}{\mathbf{SDFrm}}
\title{Grothendieck Topologies Are Extensional Presentations of the Form of Sieves}
\author{Roy Ferguson, Zurab Janelidze}
\date{Source: arXiv:2609.14671v1 (CC BY 4.0)}
\begin{document}
\maketitle

\noindent\emph{Source and licence.} Grothendieck Topologies Are Extensional Presentations of the Form of Sieves. Version arXiv:2609.14671v1 (CC BY 4.0), distributed by arXiv under the Creative Commons Attribution 4.0 International licence, http://creativecommons.org/licenses/by/4.0/, as recorded in the arXiv licence metadata for this exact version and shown on the abstract page https://arxiv.org/abs/2609.14671v1. Full licence notice: \texttt{licenses/arxiv-2609.14671-CC-BY-4.0.txt}.\par\smallskip

\noindent\textit{Editorial scope.} Counted unit: the article's theorem thm:cokernels-not-pullback-stable (source0.tex lines 1216--1220). The article's complete original proof of it is delivered with it (source0.tex lines 1222--1319, one \texttt{proof} environment of 98 lines). No mathematics is written, repaired or re-derived: the statement and the proof are the article's own lines, in the article's own notation, and no retained line is substituted or re-flowed.  Retained uncounted as the source-local context the proof needs: lines 831--840; lines 1152--1156.  Nothing was omitted from any retained span, so the package carries no omission record.  No retained line contains a citation, so the wrapper carries no bibliography and no bibliography entry.  No retained line contains a figure environment, an image-inclusion command or drawing code, so no image is delivered and the package declares no asset dependency.  Editorial formatting only, in addition: an AMS \texttt{amsart} preamble that restates the article's own declarations and macros used by the retained lines, namely the environments \texttt{theorem} on separate counters, together with the standard packages those lines need.  The retained environments print in this document as Theorem 1 in order of appearance, whereas the article prints the same items with the numbering of its own full document.  The complete native source of the article as distributed in the arXiv e-print for this exact version is bundled unchanged as \texttt{sources/210345/source0.tex}.
\begin{theorem}\label{thm:limits-and-monomorphisms}
Subject to the usual size conventions, both \(\DFrm(\C)\) and
\(\SDFrm(\C)\) have all small limits.  These limits are created by the
forgetful functors to indexed meet-semilattices.  Moreover, a morphism
\(u:F\to G\) is a monomorphism if and only if every component
\[
                       u_X:F_X\longrightarrow G_X
\]
is injective.  In that case every \(u_X\) is an order embedding.
\end{theorem}

\begin{theorem}\label{thm:composition-normal-maps}
Kernels need not be closed under composition in either
\((\DFrm(\C),\Null)\) or \((\SDFrm(\C),\Null_{\mathrm s})\).  Cokernels,
by contrast, are closed under composition in both categories.
\end{theorem}

\begin{theorem}\label{thm:cokernels-not-pullback-stable}
Relative cokernels are not stable under pullback in either
\(\DFrm(\C)\) or \(\SDFrm(\C)\).  The failure occurs already when
\(\C=\mathbf1\), and even for pullback along a top-reflecting morphism.
\end{theorem}

\begin{proof}
Put
\[
 S=\{r,p\}\amalg\mathbb N,\qquad
 T=\{p\}\amalg\mathbb N,\qquad
 I=\{*\}\amalg\mathbb N,
\]
and let \(M=\mathcal P(S)\), with intersection as meet.  Set
\[
 A_* = B_*=\{p\},\qquad A_n=\{n\},\qquad B_n=\varnothing .
\]
Equip \(M\) with the distributivity generated by the two exact
presentations
\begin{equation}\label{eq:pullback-counterexample-presentations}
                    T=\bigvee_{i\in I}A_i,
             \qquad \{p\}=\bigvee_{i\in I}B_i .
\end{equation}
Let
\[
 D=\{d\subseteq S\mid r,p\in d\text{ and }
                     \mathbb N\setminus d\text{ is finite}\}.
\]
This filter is normal.  Indeed, its characteristic meet morphism
\(\chi_D:M\to\mathbf2\), defined by \(\chi_D(C)=1\) precisely when
\(C\in D\), becomes a distributivity morphism if \(\mathbf2\) is
equipped with the distributivity generated by
\[
                          0=\bigvee_{i\in I}0.
\]
Both presentations in
\eqref{eq:pullback-counterexample-presentations} map to this
presentation.  Thus \(D=\ker\chi_D\).

Let
\[
                    q:M\longrightarrow Q=M/\theta_D
\]
be the cokernel of \(D\hookrightarrow M\).  For each \(n\), the member
\(d_n=S\setminus\{n\}\) lies in \(D\).  Since
\(d_n\mathrel{\theta_D}S\), meeting with \(\{n\}\) gives
\[
                     \{n\}\mathrel{\theta_D}\varnothing .
\]
Compatibility with the two selected presentations therefore yields
\[
                     q(T)=q(\{p\})=:a.
\]
Moreover \(a\neq\top\), because the top class of \(\theta_D\) is
exactly \(D\), whereas neither \(T\) nor \(\{p\}\) belongs to \(D\).

Give the two-element meet-semilattice \(H=\{0<1\}\) its least
distributivity and define
\[
                   j:H\longrightarrow Q,\qquad
                   j(0)=a,\quad j(1)=\top .
\]
This is a top-reflecting distributivity morphism.  Form the pullback
\[
\begin{array}{ccc}
P=M\times_QH&\xrightarrow{\pi_M}&M\\
{\scriptstyle\pi_H}\downarrow&&\downarrow{\scriptstyle q}\\
H&\xrightarrow{j}&Q .
\end{array}
\]
By Theorem~\ref{thm:limits-and-monomorphisms}, the distributivity of
\(P\) is detected coordinatewise.  Since \(H\) has only unit
presentations, so does \(P\).  The two elements
\[
                         x=(T,0),\qquad y=(\{p\},0)
\]
belong to \(P\) and satisfy \(\pi_H(x)=\pi_H(y)\).

The kernel of \(\pi_H\) is the filter
\[
                   \widetilde D=\{(d,1)\mid d\in D\}.
\]
For a meet-semilattice with only unit presentations, the least
congruence killing a filter is the common-witness congruence
\[
 z\equiv_{\widetilde D}z'
 \quad\Longleftrightarrow\quad
 z\wedge k=z'\wedge k
 \quad\text{for some }k\in\widetilde D .
\]
If this relation identified \(x\) and \(y\), some \(d\in D\) would
satisfy
\[
                         T\cap d=\{p\}\cap d=\{p\}.
\]
This is impossible, since \(d\cap\mathbb N\) is cofinite.  Hence the
cokernel of \(\ker\pi_H\) does not identify \(x\) and \(y\), whereas
\(\pi_H\) does.  Every cokernel is the cokernel of its own kernel, as
shown in Theorem~\ref{thm:composition-normal-maps}; therefore
\(\pi_H\) is not a cokernel.

Over the terminal base every distributivity is strong.  The same
example proves the assertion in both categories.
\end{proof}

\end{document}
